\chapter{Getting Access: Rates and Credit}\label{ch:m2:getting-access-rates-credit}

A swap is agreed in a second on a screen. Before the first one, a fund needs a master
agreement with each dealer, negotiated schedule by schedule, a \omterm{def:m2:swap-clearing-and-the-ccp-basis:umr}{credit support annex} that
says what collateral moves and when, an identifier that tells regulators who it is, and
accounts with a \omterm{def:m2:getting-access-rates-credit:cb}{clearing broker} for the swaps that must be cleared. Treasury repo needs
its own master agreement and, increasingly, a way into the clearing house: the United
States requires eligible repo to be cleared from June 2027. For each piece the fund
chooses between renting someone else's balance sheet and membership, and building its
own. This chapter explains the documents, the status a firm has on the trading
platforms, how non-banks reach the inter-dealer markets, \omterm{def:m2:repo-and-specials:sponsored}{sponsored repo} and clearing
membership, and swap \omterm{def:m2:getting-access-rates-credit:cb}{clearing brokers}, and builds a model of what each route costs.

\section{Documentation before the first trade}

\begin{definition}[ISDA master agreement, Global Master Repurchase Agreement]\label{def:m2:getting-access-rates-credit:isda}
An \emph{ISDA master agreement}\index{ISDA master agreement} is the standard contract,
published by the International Swaps and Derivatives Association, that governs all the
over-the-counter derivatives two parties trade with each other: a printed form, a
schedule negotiated between them, and a confirmation for each trade, all forming one
agreement. The \emph{Global Master Repurchase Agreement}\index{Global Master Repurchase
Agreement} (GMRA) is the corresponding standard for repo, published by the
International Capital Market Association.
\end{definition}

The \omterm{def:m2:getting-access-rates-credit:isda}{ISDA master agreement} exists in 1992 and 2002 versions; its central provision is
close-out netting: on a default, every outstanding trade between the two parties is
replaced by one early termination amount owed by one to the other. That is what lets a
dealer's credit department and its regulators measure exposure net rather than trade by
trade, as long as the netting is enforceable under the insolvency law of the
counterparty's country. The \omterm{def:m2:swap-clearing-and-the-ccp-basis:umr}{credit support annex} of chapter 10 sits on top of it and
sets the collateral. The GMRA, first published in 1992 and last revised in 2011, does
the same for repo, with legal opinions on its enforceability in many jurisdictions
published by its sponsor (\cref{fig:m2:getting-access-rates-credit:docs}).

\begin{definition}[Legal entity identifier]\label{def:m2:getting-access-rates-credit:lei}
A \emph{legal entity identifier}\index{legal entity identifier} (LEI) is a unique
20-character code, under the ISO~17442 standard, that identifies a legal entity in
trade reports and regulatory filings, with public reference data on who the entity is
and who owns it.
\end{definition}

The LEI came out of the 2008 crisis, when regulators could not tell who owed what to
whom: the G20 endorsed the system's charter in 2012, and a Swiss foundation, GLEIF,
administers it. Regulators use the LEI to identify the parties in trade reports, so a
fund obtains one before it trades.

\begin{omfigure}
\centering
\begin{tikzpicture}[font=\small,
  doc/.style={draw,thick,rounded corners=2pt,align=center,minimum width=3.4cm,minimum height=0.75cm,inner sep=2pt}]
\node[doc,draw=omDef,fill=omDef!10] (ma) at (0,0) {ISDA master agreement};
\node[doc,draw=omDef,fill=omDef!5] (sch) at (0,0.95) {schedule (negotiated)};
\node[doc,draw=omThm,fill=omThm!8] (csa) at (0,1.9) {credit support annex};
\node[doc,draw=gray,fill=gray!8] (conf) at (0,2.85) {confirmations, per trade};
\node[font=\footnotesize\bfseries] at (0,3.7) {swaps, per dealer};
\node[doc,draw=omProp,fill=omProp!10] (gm) at (5,0) {GMRA 2011};
\node[doc,draw=omProp,fill=omProp!5] (ann) at (5,0.95) {annexes (negotiated)};
\node[doc,draw=gray,fill=gray!8] (rc) at (5,1.9) {confirmations, per trade};
\node[font=\footnotesize\bfseries] at (5,3.7) {repo, per counterparty};
\node[doc,draw=omThm,fill=omThm!8,minimum width=3.8cm] (clr) at (10,0) {clearing agreements};
\node[doc,draw=gray,fill=gray!8,minimum width=3.8cm] (plat) at (10,0.95) {platform onboarding};
\node[doc,draw=omDef,fill=omDef!8,minimum width=3.8cm] (lei) at (10,1.9) {legal entity identifier};
\node[font=\footnotesize\bfseries] at (10,3.7) {once, for everything};
\end{tikzpicture}

\omcaption{The documents a fund needs before its first swap and its first repo: an \omterm{def:m2:getting-access-rates-credit:isda}{ISDA
master agreement} with each dealer, with its schedule, \omterm{def:m2:swap-clearing-and-the-ccp-basis:umr}{credit support annex} and
confirmations; a GMRA with each repo counterparty; and, once, clearing agreements,
platform onboarding and an LEI. Schematic.}\label{fig:m2:getting-access-rates-credit:docs}
\end{omfigure}

\section{Dealer and client status on the platforms}

Rates and credit trade on two kinds of electronic platform (chapters 4 and 22): the
inter-dealer markets, where dealers and a few large non-banks trade with each other,
mostly on order books; and the \omterm{def:m2:the-treasury-market:idb}{dealer-to-client platforms}, where clients request quotes
from dealers. A firm's status decides what it sees and whom it can trade with: as a
client, it asks for prices; as a dealer, it answers requests and may see flows that
clients do not.

In swaps, the rules push against closed clubs: an American \omterm{def:m2:interest-rate-swaps:sef}{swap execution facility} must
give any eligible contract participant impartial access to its markets, under criteria
that are impartial, transparent and applied in a fair and non-discriminatory manner,
with comparable fees for comparable access. Access to a platform is still not access
to trading: the counterparties on it must accept the firm's credit, directly or
through its \omterm{def:m2:getting-access-rates-credit:cb}{clearing broker}.

\section{Inter-dealer access for non-banks}

The inter-dealer Treasury market is no longer a dealers' club. Chapter 4 showed that, on
15 October 2014, principal trading firms, non-banks trading for their own account,
accounted for more than half of the volume in the inter-dealer cash and futures
markets studied. They reach the order books through the platforms' own onboarding and,
for credit and settlement, through a bank or a clearing member that stands behind their
trades, the same pattern as the \omterm{def:m2:the-fx-market:sdp}{FX prime brokerage} of the last chapter. For a
non-bank, the order book gives anonymous access to the tightest prices; the price of
it is the credit arrangement, and the obligation to settle and clear what it trades.

\section{Sponsored repo and clearing membership}

\begin{definition}[Sponsored member]\label{def:m2:getting-access-rates-credit:sponsored}
A \emph{sponsored member}\index{sponsored member} is a firm admitted to a clearing
house through a sponsoring member, which submits the sponsored member's trades for
clearing and guarantees its obligations to the clearing house.
\end{definition}

Chapter 5 described \omterm{def:m2:repo-and-specials:sponsored}{sponsored repo}. The Fixed Income Clearing Corporation, the clearing
house for Treasury repo, offers three routes to a fund
(\cref{fig:m2:getting-access-rates-credit:routes}). It can become a \omterm{def:m2:getting-access-rates-credit:sponsored}{sponsored member} of
a bank that is a sponsoring member: in November 2024 the SEC approved rule changes that
removed the requirement that \omterm{def:m2:getting-access-rates-credit:sponsored}{sponsored members} be qualified institutional buyers, and
bank sponsors must have equity capital of at least USD~5 billion and be
well capitalised. It can use a member's agent clearing service, which FICC now presents
as a ``done-away'' model: the fund trades with anyone and the agent clears. Or it can
apply for membership itself, with the capital, operations and contributions to the
clearing fund that brings. A fourth route, uncleared \omterm{def:m2:repo-and-specials:triparty}{bilateral repo}, narrows as the
\omterm{def:m2:swap-clearing-and-the-ccp-basis:mandate}{clearing mandate} takes effect.

\begin{omfigure}
\centering
\begin{tikzpicture}[font=\small,
  box/.style={draw,thick,rounded corners=2pt,align=center,minimum height=0.95cm,minimum width=2.8cm,inner sep=3pt},
  arr/.style={-{Stealth[length=2.5mm]},thick}]
\node[box,draw=omDef,fill=omDef!8] (f) at (0,0) {fund};
\node[box,draw=omThm,fill=omThm!8] (s) at (5,1.6) {sponsoring bank};
\node[box,draw=omProp,fill=omProp!8] (a) at (5,0) {agent clearing member};
\node[box,draw=gray,fill=gray!15] (c) at (10.2,0) {FICC\\(clearing house)};
\node[box,draw=gray,fill=gray!8] (d) at (5,-1.6) {dealer (bilateral)};
\draw[arr] (f) -- node[above,sloped,font=\footnotesize] {rent: sponsored} (s);
\draw[arr] (s) -- (c);
\draw[arr] (f) -- node[above,font=\footnotesize] {done-away} (a);
\draw[arr] (a) -- (c);
\draw[arr,dashed] (f) -- node[below,sloped,font=\footnotesize] {uncleared} (d);
\draw[arr,very thick,omDef] (f.north) to[out=60,in=150] node[pos=0.5,above=2pt,font=\footnotesize] {buy: direct membership} (c.north);
\end{tikzpicture}

\omcaption{Routes to cleared Treasury repo for a fund: sponsored membership through a
bank, an agent clearing member that clears trades done with others, direct membership of
the clearing house, and uncleared \omterm{def:m2:repo-and-specials:triparty}{bilateral repo} with a dealer, which the \omterm{def:m2:swap-clearing-and-the-ccp-basis:mandate}{clearing
mandate} narrows. Schematic.}\label{fig:m2:getting-access-rates-credit:routes}
\end{omfigure}

\begin{dated}{2026-09}{Treasury clearing and sponsored access}\label{dat:m2:getting-access-rates-credit:access}
SEC Treasury clearing rule: compliance dates 31 December 2026 for eligible cash
transactions and 30 June 2027 for eligible repo (chapter 4). FICC's Sponsored Service:
average daily volume of USD~2.3 trillion in the second quarter of 2026, by DTCC. SEC
approval of FICC's access rule changes, including the removal of the qualified
institutional buyer requirement for \omterm{def:m2:getting-access-rates-credit:sponsored}{sponsored members}: 21 November 2024.
\end{dated}

\section{Swap clearing brokers}

\begin{definition}[Clearing broker (swaps)]\label{def:m2:getting-access-rates-credit:cb}
A \emph{clearing broker}\index{clearing broker (swaps)} for swaps is a clearing member
of a swap clearing house that clears its clients' swaps, posts their margin to the
clearing house, and guarantees their performance to it; in the United States it is a
futures commission merchant registered for swaps.
\end{definition}

Swaps subject to the \omterm{def:m2:swap-clearing-and-the-ccp-basis:mandate}{clearing mandate} of chapter 10 must be cleared, so a fund needs a
\omterm{def:m2:getting-access-rates-credit:cb}{clearing broker} at the clearing houses where its dealers want to face it. The broker
charges fees on the notional or per ticket, requires initial margin, often with an
add-on for the fund's risk, and limits the fund's positions by its own credit
appetite. Becoming a clearing member directly is possible for large firms, but brings
default-fund contributions, default-management duties and operational requirements
that only large books justify.

\begin{proposition}[When to buy access]\label{prop:m2:getting-access-rates-credit:breakeven}
Let a route have a fixed annual cost $F$ and a variable annual cost $u$ per unit of
repo balance, for a book whose swap notional is a fixed multiple of its repo balance. If
building access costs $F_b > F_r$ but $u_b < u_r$ against renting it, building is
cheaper exactly when the repo balance exceeds
\[
B^* = \frac{F_b - F_r}{u_r - u_b}.
\]
\end{proposition}

\begin{proof}
The annual costs are $F_r + u_r B$ and $F_b + u_b B$; the second is smaller when $(u_r -
u_b)B > F_b - F_r$.
\end{proof}

\begin{example}[Three routes]\label{ex:m2:getting-access-rates-credit:routes}
In the chapter's illustrative model, renting costs USD~0.5 million a year fixed and 11.5
basis points a year per dollar of repo balance, including the swaps and funded margin;
buying costs USD~8 million fixed and 5.9 basis points, including the capital in the
clearing fund; bilateral costs 0.2 million and 18.5 basis points. A fund with USD~5
billion of repo and 10 billion of swaps pays 6.25 million a year renting, 10.95 million
buying and 9.45 million bilaterally. Bilateral beats renting only below USD~0.43
billion of repo; buying beats renting above USD~13.4 billion
(\cref{fig:m2:getting-access-rates-credit:costs}).
\end{example}

\begin{omfigure}
\begin{tikzpicture}
\begin{axis}[width=11cm,height=5cm,
  xlabel={repo balance (USD bn), swaps twice as large},ylabel={annual cost (USD m)},
  tick label style={font=\small},label style={font=\small},
  xmin=0,xmax=30,ymin=0,ymax=40,grid=major,grid style={gray!25},
  legend columns=3,legend style={font=\footnotesize,at={(0.5,-0.3)},anchor=north,draw=none,fill=none,/tikz/every even column/.append style={column sep=8pt}},
  table/col sep=comma]
\addplot[omDef,very thick] table[x=repo,y=rent]{figdata/markets-2/28-getting-access-rates-credit/costs.csv};
\addplot[omThm,very thick] table[x=repo,y=buy]{figdata/markets-2/28-getting-access-rates-credit/costs.csv};
\addplot[omProp,very thick,dashed] table[x=repo,y=bilateral]{figdata/markets-2/28-getting-access-rates-credit/costs.csv};
\addplot[only marks,mark=*,mark size=1.8pt,black,forget plot] coordinates {(13.39,15.90)};
\node[font=\footnotesize,anchor=north west] at (axis cs:13.8,15.2) {break-even 13.4};
\legend{{rent (sponsored, broker)},buy (direct),bilateral}
\end{axis}
\end{tikzpicture}

\omcaption{Annual cost of access to Treasury repo and swaps against the repo balance,
with swap notional twice as large, for three routes: renting (\omterm{def:m2:repo-and-specials:sponsored}{sponsored repo} and a swap
\omterm{def:m2:getting-access-rates-credit:cb}{clearing broker}), buying (direct memberships) and bilateral. Buying pays above USD~13.4
billion of repo. Illustrative costs; data: the chapter's tutorial.}\label{fig:m2:getting-access-rates-credit:costs}
\end{omfigure}

The break-even is only as good as its inputs, and the most uncertain is the spread the
sponsor charges, which depends on its balance sheet and its appetite
(\cref{fig:m2:getting-access-rates-credit:sens}).

\begin{omfigure}
\begin{tikzpicture}
\begin{axis}[width=11cm,height=4.6cm,
  xlabel={sponsor's spread on the repo balance (bp a year)},ylabel={break-even (USD bn)},
  tick label style={font=\small},label style={font=\small},
  xmin=4,xmax=12,ymin=0,ymax=50,grid=major,grid style={gray!25},
  table/col sep=comma]
\addplot[omDef,very thick] table[x=spread,y=breakeven]{figdata/markets-2/28-getting-access-rates-credit/sensitivity.csv};
\addplot[only marks,mark=*,mark size=1.8pt,black] coordinates {(5,28.85) (8,13.39) (12,7.81)};
\end{axis}
\end{tikzpicture}

\omcaption{The repo balance above which direct membership is cheaper than sponsorship,
against the sponsor's spread, other costs unchanged. At 5 basis points it is USD~28.8
billion, at 8 basis points 13.4 and at 12 basis points 7.8. Illustrative; data: the
chapter's tutorial.}\label{fig:m2:getting-access-rates-credit:sens}
\end{omfigure}

\section{Tutorial: three ways to access repo and swaps}\label{tut:m2:getting-access-rates-credit:cost}

\begin{tutorial}
\textbf{Goal.} Cost three routes to Treasury repo and swaps for a mid-size fund, find the
break-even sizes, and see how they move with the sponsor's spread and the fixed cost of
membership.
\textbf{End state:} \cref{fig:m2:getting-access-rates-credit:costs,fig:m2:getting-access-rates-credit:sens},
\cref{ex:m2:getting-access-rates-credit:routes} and the numbers of the weekend problem.
\begin{tutsteps}
\item \textbf{The model}: routes, annual costs, unit costs, break-evens.
\omcode{code/firm/clearcost/firm_clearcost.py}{12}{51}{The access-cost model.}\label{lst:m2:getting-access-rates-credit:model}
\item \textbf{Run} \texttt{access\_cost\_demo.tutorial()},
\texttt{access\_cost\_demo.problem()} and \texttt{fig\_access\_cost.py}.
\end{tutsteps}
\textbf{What to change next.} Make the sponsor's spread rise with the fund's balance, as
a bank's balance-sheet cost would, and find whether the break-even still exists; then
add the value of the netting a direct member gets across all its counterparties.
\end{tutorial}

\section{Build: the access-cost model}\label{bld:m2:getting-access-rates-credit:clearcost}

\begin{build}
\bfield{Purpose} The miniature firm decides how it reaches each market, and revisits
the decision as it grows.
\bfield{Interface} \texttt{Route(name, fixed, repo\_bp, swap\_bp, repo\_margin,
swap\_margin, fund\_share, capital\_rate)}; \texttt{annual\_cost(route, repo, swaps,
funding)}; \texttt{variable\_cost}; \texttt{unit\_cost(route, swaps\_per\_repo,
funding)}; \texttt{breakeven(a, b, swaps\_per\_repo, funding)}; \texttt{cheapest(routes,
repo, swaps, funding)}.
\bfield{Rules} Costs a year in dollars; fees in basis points of balance or notional;
margin funded at a spread; clearing-fund capital at a cost of capital; costs linear in
the book's size beyond the fixed part.
\bfield{Acceptance tests} \texttt{code/firm/clearcost/tests/}: the cost is fixed plus
linear; at the break-even the two routes cost the same, below it renting is cheaper and
above it buying is; a dominated route has no break-even.
\bfield{Stretch} Costs that rise with size; the value of netting and of anonymity;
default-fund and loss-allocation exposures of a direct member; stress costs when a
sponsor withdraws (One Quant Book 16 on the firm).
\end{build}

\begin{omsources}
\item ISDA, ``Legal guidelines for smart derivatives contracts: the ISDA Master
Agreement'', February 2019.
\item ICMA, Global Master Repurchase Agreement; GLEIF, introducing the legal entity
identifier.
\item SEC, Release No.~34-101694 (21 November 2024) approving FICC's access rule changes;
DTCC, FICC Sponsored Service.
\item 17 CFR 37.202, access requirements for swap execution facilities.
\end{omsources}

\section{Exercises}

\begin{exercise}[$\star$]\label{exo:m2:getting-access-rates-credit:1}
What does close-out netting do, and why does a dealer care whether it is enforceable in
its client's country?
\end{exercise}

\begin{exercise}[$\star$]\label{exo:m2:getting-access-rates-credit:2}
List the documents a fund needs before its first cleared swap.
\end{exercise}

\begin{exercise}[$\star$]\label{exo:m2:getting-access-rates-credit:3}
What does an LEI identify, and what two questions does its reference data answer?
\end{exercise}

\begin{exercise}[$\star\star$]\label{exo:m2:getting-access-rates-credit:4}
Rent costs 0.5 million a year and 11.5 basis points; buy costs 8 million and 5.9 basis
points. Compute the break-even from the formula.
\end{exercise}

\begin{exercise}[$\star\star$]\label{exo:m2:getting-access-rates-credit:5}
What is the difference between sponsored clearing and a done-away agent clearing
service, from the fund's point of view?
\end{exercise}

\begin{exercise}[$\star\star$]\label{exo:m2:getting-access-rates-credit:6}
Why does the Treasury \omterm{def:m2:swap-clearing-and-the-ccp-basis:mandate}{clearing mandate} change the comparison between \omterm{def:m2:repo-and-specials:triparty}{bilateral repo} and
the cleared routes?
\end{exercise}

\begin{exercise}[$\star\star\star$]\label{exo:m2:getting-access-rates-credit:7}
\emph{Coding.} Remove the swaps (a repo-only fund) and find the break-even between
renting and buying. Then raise the fixed cost of membership to USD~12 million, swaps
back in. Explain both changes.
\end{exercise}

\begin{exercise}[$\star\star\star$]\label{exo:m2:getting-access-rates-credit:8}
\emph{Find the flaw.} ``Direct membership is always cheaper in the long run: the fees
are lower, and fixed costs are paid only once.'' Correct it.
\end{exercise}

\section{Problem: Build or Rent}

\begin{problem}[{Weekend problem --- when direct membership pays}]\label{pb:m2:getting-access-rates-credit:1}
A fund's rates book has a repo balance $B$ and a swap notional $2B$. Renting access
(\omterm{def:m2:repo-and-specials:sponsored}{sponsored repo}, a swap \omterm{def:m2:getting-access-rates-credit:cb}{clearing broker}) costs USD~0.5 million a year, 8 basis points a
year on the repo balance, 0.5 on swap notional, with margin of 2\% of the repo and 1.5\% of
the swaps funded at 0.5\% a year. Buying it (direct memberships) costs USD~8 million a
year, 1 and 0.2 basis points, the same margin, and a clearing-fund contribution of 0.2\%
of the repo balance that costs 10\% a year in capital. Bilateral, uncleared, costs 0.2
million, 12 and 1 basis points, with margin of 3\% on both. All figures are illustrative.

\medskip\noindent\textbf{Part I --- Unit costs.}
\begin{enumerate}
\item Compute each route's variable cost per dollar of repo balance, in basis points.
\item What does each route cost for USD~5 billion of repo and 10 billion of swaps?
\item Which is cheapest at that size?
\item Below which size is bilateral cheapest?
\item Why does the clearing-fund contribution count as a cost?
\end{enumerate}

\medskip\noindent\textbf{Part II --- The break-even.}
\begin{enumerate}[resume]
\item Find the repo balance above which buying beats renting.
\item What does either route cost at that size?
\item What would buying save at USD~20 billion of repo?
\item How does the break-even change for a repo-only fund?
\item How does it change if membership costs 12 million a year?
\end{enumerate}

\medskip\noindent\textbf{Part III --- Sensitivity.}
\begin{enumerate}[resume]
\item What is the break-even if the sponsor charges 5 basis points? 12?
\item Why is the sponsor's spread the most uncertain input?
\item What costs of direct membership does the model leave out?
\item What does renting cost that the model does not show?
\item How does the \omterm{def:m2:swap-clearing-and-the-ccp-basis:mandate}{clearing mandate} change the bilateral route?
\end{enumerate}

\medskip\noindent\textbf{Part IV --- Judgement.}
\begin{enumerate}[resume]
\item The fund expects to grow from 5 to 20 billion over three years. When should it
start building?
\item Why might a fund rent even above the break-even?
\item What documents must be in place whichever route it chooses?
\item State the \emph{named result}: the balance-sheet size at which direct membership
beats a \omterm{def:m2:getting-access-rates-credit:cb}{clearing broker}, and its cost.
\item In one sentence: what does a fund buy when it becomes a member?
\end{enumerate}
\end{problem}

\section{Interview questions}

\begin{interviewq}[$\star$ \iqroles{trader, developer}]\label{iq:m2:getting-access-rates-credit:1}
What is an \omterm{def:m2:getting-access-rates-credit:isda}{ISDA master agreement}, and why does close-out netting matter?
\end{interviewq}

\begin{interviewq}[$\star$ \iqroles{trader, bank}]\label{iq:m2:getting-access-rates-credit:2}
How does a hedge fund get access to cleared Treasury repo?
\end{interviewq}

\begin{interviewq}[$\star\star$ \iqroles{researcher, risk}]\label{iq:m2:getting-access-rates-credit:3}
Build a model to decide between a \omterm{def:m2:getting-access-rates-credit:cb}{clearing broker} and direct membership.
\end{interviewq}

\begin{interviewq}[$\star\star$ \iqroles{trader}]\label{iq:m2:getting-access-rates-credit:4}
What does a swap \omterm{def:m2:getting-access-rates-credit:cb}{clearing broker} do, and what can it do to you in a stress?
\end{interviewq}

\begin{interviewq}[$\star\star$ \iqroles{trader, researcher}]\label{iq:m2:getting-access-rates-credit:5}
How do non-banks reach the inter-dealer Treasury market?
\end{interviewq}

\begin{interviewq}[$\star\star\star$ \iqroles{developer}]\label{iq:m2:getting-access-rates-credit:6}
Design the onboarding system that tracks, for each counterparty, the documents signed,
their key terms and what they allow the firm to trade.
\end{interviewq}
